\documentclass{article}
\usepackage[utf8]{inputenc}
\usepackage[T1]{fontenc}
\usepackage{geometry}
\geometry{
    a4paper,
    left=2.5cm,
    right=2.5cm,
    top=2.5cm,
    bottom=2.5cm
}
\usepackage{graphicx}
\usepackage{float}
\usepackage{hyperref}
\usepackage[english]{babel}
\usepackage{amsmath}
\usepackage{amssymb}
\usepackage{array}
\usepackage{booktabs}
\usepackage{caption}
\usepackage{subcaption}
\usepackage{enumitem}
\usepackage{fancyhdr}
\usepackage{titlesec}
\usepackage{url}
\usepackage{textcomp}
\usepackage{xcolor}
\usepackage{listings}
\usepackage{multirow}
\usepackage{longtable}
\usepackage{tabularx}
\usepackage{times}

% Set up page style
\pagestyle{fancy}
\fancyhf{}
\rfoot{\thepage}
\renewcommand{\headrulewidth}{0pt}

% Title formatting
\titleformat{\chapter}[display]{\normalfont\bfseries}{\chaptertitle\thechapter}{0pt}{\Huge}
\titlespacing*{\chapter}{0pt}{50pt}{40pt}

% Document info
\title{Template Document}
\author{Academic Support Agent}
\date{\today}

\begin{document}
\maketitle

%Begin Inicia o Documento



% Content from cabecalho would be included here
 %Alterar T tulo, autores, e apagar tabela !

%----------------------------------------TEXTOS-------------------------------------------------%


\vspace{0.5cm}
\noindent\textbf{RESUMO}: O prop sito destas instru es   orientar aos autor(es) quanto   formata o dos resumos expandidos a serem submetidos ao Congresso de Inova o, Ci ncia e Tecnologia do Instituto Federal de Educa o, Ci ncia e Tecnologia de S o Paulo. Os documentos devem ser redigidos de acordo com as normas para elabora o do resumo expandido. O arquivo de submiss o dever  estar desbloqueado no formato \textit{portable document format} (pdf) compat vel com o Adobe Acrobat Reader(TM). O texto deve iniciar na mesma linha do item, ser claro, sucinto e, obrigatoriamente, explicar o(s) objetivo(s) pretendido(s), procurando justificar sua import ncia (sem incluir refer ncias bibliogr ficas), os principais procedimentos adotados, os resultados mais expressivos e conclus es, contendo no m ximo 200 palavras. N o dever  conter f rmulas e cita es e refer ncias bibliogr ficas. O resumo expandido apresentado no evento ser  publicado nos Anais (\textit{ISSN}: 2178-9959). O texto com as instru es e em par nteses devem ser removidos do documento final. (Times New Roman, 11, Justificado, M ximo 200 palavras).

\vspace{0.5cm}
\noindent\textbf{PALAVRAS-CHAVE}: m ximo de seis, separadas por ponto e v rgula (;), procurando n o repetir palavras do t tulo, escritas em letras min sculas. (Times New Roman, 11, Justificado).

\vspace{0.5cm}
 \begin{center}
 \textbf{T TULO EM INGL S (Times New Roman, 11, Negrito, Centralizado)}
 \end{center}

\noindent\textbf{ABSTRACT}: Tradu o do resumo para a l ngua inglesa. (Times New Roman, 11, Justificado).

\vspace{0.5cm}
\noindent\textbf{KEYWORDS}: Tradu o das palavras-chave para a l ngua inglesa. (Times New Roman, 11, Justificado).

\section*{Introdu o}

No m ximo 20 linhas, evitar divaga es, utilizando bibliografia apropriada para formular os problemas abordados e a justificativa da import ncia do assunto, deixando claro a(s) hip tese(s) e o(s) objetivo(s) do trabalho. (Times New Roman, 11, Justificado, M ximo 20 linhas). De acordo com \cite{alves} podemos...

\section*{Materiais e M todos}

Os materiais e m todos utilizados no desenvolvimento da pesquisa devem ser adequadamente descritos. (Times New Roman, 11, Justificado).

\textbf{Modelo de Equa o:}

\begin{equation} \label{eq:nome}
IC=\frac{F*9,81}{A}*10^{-6}
\end{equation}
em que,

IC -  ndice de cone,$MPa$;

F - for a, $kgf$;

A -  rea do cone,$m^{2}$.

%De acordo com a \cite{carvalho} 

\section*{Resultados e Discuss o}

Ilustra es e gr ficos devem ser apresentados com tamanho e detalhes suficientes para a composi o gr fica final, preferivelmente na mesma posi o do texto.

\textbf{Gr ficos:} devem apresentar-se sem bordas, descritos com o mesmo tipo e tamanho de letras contidas no texto e a legenda na posi o inferior do mesmo. A numera o deve ser sucessiva em algarismos ar bicos.

\textbf{Tabelas:} evitar tabelas extensas e dados sup rfluos; adequar seus tamanhos ao espa o  til do papel e colocar, na medida do poss vel, apenas linhas cont nuas horizontais; suas legendas devem ser concisas e autoexplicativas. Na discuss o, confrontar os dados obtidos com a literatura. (Times New Roman, 11, Justificado).
\vspace{0.5cm}

\noindent\textbf{Modelos de Figuras:}


\begin{figure}[ht]
\centering
\includegraphics[width=10cm,angle=0]{grafico_1}
\caption{Mapas de teor de  gua das camadas de 0-2,2 e 0,2-0,4 m de profundidade.}
\label{fig:nome_referencia_figura1}
\end{figure}

\begin{figure}[ht]
\centering
\includegraphics[width=10cm,angle=0]{grafico_2.png}
\caption{Mapas do  ndice de cone (MPa) referente aos dados coletados nas diferentes profundidades nas linhas e nas entrelinhas da cultura da cana.}
\label{fig:nome_referencia_figura2}
\end{figure}

\newpage
\noindent\textbf{Modelo de Tabela:}

\begin{table}[ht]
\caption{An lise do IC nas linhas (L) e entrelinhas (E) de cana nas diferentes profundidades amostradas pelo  ndice de cone.}
\vspace{.2cm}
\centering
\resizebox{\textwidth}{!}{%
\begin{tabular}{ccccccccc}
\hline
Profundidade (m) & \multicolumn{2}{c}{0 a 0,1} & \multicolumn{2}{c}{0,1 a 0,2} & \multicolumn{2}{c}{0,2 a 0,3} & \multicolumn{2}{c}{0,3 a 0,4} \\ \hline
                 & L            & E            & L             & E             & L             & E             & L             & E             \\
M dia (MPa)      & 1,39**       & 4,28**       & 1,86**        & 4,29**        & 2,20**        & 3,83**        & 2,46**        & 3,44**        \\
CV(\%)           & 54           & 57           & 55            & 54            & 46            & 49            & 48            & 43            \\ \hline
\end{tabular}%
}
\small{**:valores significativos para o n vel de signific ncia de 1\% pelo teste de Tukey; L -- linhas; E -- entrelinhas.}
\label{tab:nome_tabela1}
\end{table}

\section*{Conclus es}

Devem basear-se exclusivamente nos resultados do trabalho. Evitar a repeti o dos resultados em listagem subsequente, buscando, sim, confrontar o que se obteve com os objetivos inicialmente estabelecidos. (Times New Roman, 11, Justificado).



\section*{Agradecimentos}

Inserir ap s as conclus es, de maneira sucinta. Se o projeto for financiado por alguma ag ncia de fomento, citar a fonte. (Times New Roman, 11, Justificado).



%Refer ncias Bibliogr ficas
%\textbf{REFER NCIAS}

As refer ncias devem ser listadas em ordem alfab tica \cite{galvani}, \cite{national_instruments} e \cite{pandorfi}. 



%\printbibliography[title={REFER NCIAS}]
%\printbibliography{bibliografia}
\end{document}